\chapter*{Introduction générale}
\addcontentsline{toc}{chapter}{Introduction générale}

Le crédit occupe une place centrale dans le financement de l'activité économique. En accordant un prêt, une banque ne se contente pas de mettre des fonds à la disposition d'un client : elle engage ses propres ressources sur la base d'une appréciation du risque que ce client ne puisse, en tout ou en partie, honorer ses engagements. Cette appréciation conditionne directement la solidité financière de l'établissement : un risque mal maîtrisé peut se traduire par des pertes, tandis qu'une prudence excessive prive la banque de financements par ailleurs rentables. La gestion du crédit repose ainsi sur un équilibre permanent entre plusieurs préoccupations difficiles à concilier : limiter le risque de défaut, anticiper les pertes qu'il peut engendrer, et disposer de fonds propres suffisants pour absorber les pertes inattendues.

C'est dans ce contexte que s'inscrit le projet de fin d'études présenté dans ce rapport, consacré à la gestion du risque de crédit bancaire et réalisé au sein de l'entité AI~\&~Data d'EY~Tunisie. Ce rapport constitue également le mémoire de stage de fin d'études du Master Actuariat, dont il valide la dernière étape académique. Ce domaine nécessite de faire converger l'évaluation quantitative du risque et la compréhension du fonctionnement réel d'une banque afin de concevoir des outils utiles à la décision. Le projet porte ainsi sur plusieurs dimensions complémentaires : le risque de crédit associé à chaque client, les pertes attendues, le capital nécessaire pour absorber les pertes inattendues et la rentabilité du financement accordé. Leur prise en compte constitue le fil conducteur de ce rapport. La problématique à laquelle répond le projet ainsi que ses objectifs sont précisés au chapitre 1, section 1.3.

Pour y répondre, le travail a été construit autour de deux volets complémentaires. Le premier, désigné dans ce rapport sous le nom d'\emph{Origination}, concerne l'évaluation d'une nouvelle demande de crédit, depuis l'estimation du risque jusqu'à la décision d'octroi. Le second, désigné sous le nom de \emph{Portfolio Analytics}, porte sur l'analyse d'un portefeuille déjà constitué, pour en suivre le risque, les pertes attendues, le capital mobilisé et la rentabilité dans la durée. Ces deux volets s'appuient sur un socle commun d'évaluation du risque, mais répondent à deux échelles de décision différentes : celle d'un dossier individuel, et celle d'un ensemble de crédits déjà accordés.

La réalisation du projet s'est appuyée sur une démarche itérative et incrémentale, de la compréhension du besoin métier jusqu'à la validation des traitements développés, avec des vérifications régulières permettant d'évaluer la cohérence des choix effectués et d'intégrer les ajustements nécessaires au fur et à mesure de l'avancement du travail.

Ce rapport est organisé en six chapitres. Le premier présente l'organisme d'accueil, le contexte bancaire du projet et l'étude ayant conduit à la solution proposée. Le second précise les besoins fonctionnels et non fonctionnels ainsi que les utilisateurs concernés. Le troisième décrit la préparation des données et la construction des modèles prédictifs de risque. Le quatrième explique comment ces résultats sont transformés en indicateurs financiers, puis en décisions. Le cinquième présente la conception et l'architecture de la plateforme développée. Le sixième, enfin, détaille sa réalisation et sa validation. Le rapport se referme sur une conclusion générale qui revient sur la démarche suivie et sur les apports de ce travail.
